\chapter{Il modello}
L'obiettivo è rappresentare in modo sintetico come la composizione del gruppo e lo stato di affaticamento medio varino nel tempo in risposta al carico di allenamento, agli infortuni e ai flussi di ingresso e uscita dalla popolazione. Il modello non ha la pretesa di descrivere il comportamento individuale di ogni atleta, ma di catturare le tendenze aggregate di un gruppo sufficientemente numeroso, per il quale le dinamiche collettive sono più rilevanti delle fluttuazioni individuali. In altre parole, il modello si colloca a un livello di astrazione intermedio: sufficientemente dettagliato da distinguere stati diversi degli atleti e da descrivere l'evoluzione della fatica, ma sufficientemente aggregato da rimanere trattabile analiticamente e da non richiedere la simulazione di ogni singolo individuo. \par

Il modello è formulato a tempo discreto, con $ t = 0,1,\dots,T $, dove $T$ è l'orizzonte temporale considerato. \\
Si assume che le sedute di allenamento siano distribuite circa omogeneamente nel tempo, così che ogni passo $ t \to t+1 $ corrisponda a una seduta e il tempo che intercorre tra due sedute sia approssimativamente lo stesso. Se consideriamo ad esempio un programma di allenamento basato su una seduta al giorno, possiamo dire che ogni passo $t$ corrisponda a un giorno. In questo modo è possibile confrontare anche il recupero tra due sedute consecutive, perché l'intervallo di tempo tra un passo e il successivo è pressoché costante per ipotesi. \\
Questa scelta è motivata dal fatto che gli allenamenti sono eventi discreti, separati da periodi di riposo, e che il modello intende descrivere l'evoluzione del gruppo seduta dopo seduta. L'ipotesi di distribuzione omogenea permette inoltre di trattare il tempo come una variabile regolare e di evitare la complessità aggiuntiva legata a pianificazioni irregolari, senza perdere la capacità di descrivere i fenomeni principali. \par

In questa prima formulazione, il modello non include un parametro di controllo. Le funzioni che descrivono l'efficacia della prevenzione e i carichi di allenamento sono trattate come parametri esogeni, assegnati dall'esterno. Questa scelta permette di studiare il comportamento del sistema in risposta a diversi scenari di allenamento e prevenzione, senza aggiungere la complessità dell'ottimizzazione del parametro. Pertanto l'aggiunta del controllo è rinviata a una fase successiva. In questo modo è possibile analizzare separatamente due aspetti distinti: da un lato la dinamica del sistema in condizioni assegnate, dall'altro la ricerca delle condizioni ottimali.

\section{Variabili}

Gli atleti sono suddivisi in quattro popolazioni, che rappresentano stati diversi rispetto alla condizione fisica e alla capacità di allenarsi. La suddivisione riflette la necessità di distinguere non solo tra atleti sani e infortunati, ma anche tra diversi livelli di rischio e capacità di allenarsi. Questa distinzione è importante perché il rischio di infortunio, la capacità di partecipare all'allenamento e la dinamica di recupero non sono uniformi tra tutti gli atleti, e trattarli come un unico gruppo omogeneo comporterebbe una perdita significativa di informazione. \par

Le popolazioni sono:
\begin{itemize}
    \item Atleti sani $S$: atleti non infortunati e senza infortuni pregressi. Si possono allenare senza limitazioni e sono considerati il gruppo di riferimento per il rischio di infortunio. La loro condizione fisica è ottimale e non presentano vulnerabilità residue legate a eventi precedenti.
    \item Atleti con pregressi $P$: atleti che hanno avuto infortuni prima dell'inizio del periodo preso in considerazione. Pur essendo pienamente allenabili, sono considerati più a rischio rispetto ai sani, in quanto un infortunio pregresso può lasciare una vulnerabilità residua, una ridotta stabilità articolare o una minore tolleranza al carico. Questi atleti partecipano all'allenamento senza restrizioni formali, ma la loro storia clinica giustifica un trattamento differenziato nel modello del rischio.
    \item Atleti infortunati $I$: atleti che si sono infortunati durante il periodo preso in considerazione e non si stanno allenando. Sono temporaneamente esclusi dall'attività e non contribuiscono alla densità sul tatami. La loro presenza è rilevante perché rappresenta l'esito più grave del rischio di infortunio e perché il loro recupero parziale alimenta il gruppo $M$.
    \item Atleti ``malconci'' $M$: atleti con infortuni lievi o che sono parzialmente guariti dall'infortunio e possono continuare ad allenarsi con qualche restrizione. Partecipano all'allenamento, ma la loro condizione fisica non è ottimale e il loro rischio di incorrere in un nuovo infortunio è considerato maggiore rispetto a quello degli atleti sani. Questi atleti rappresentano una categoria intermedia: non sono abbastanza compromessi da dover sospendere l'attività, ma non sono nemmeno in condizioni ottimali.
\end{itemize}

Le popolazioni $S$, $P$, $I$, $M$ sono rappresentate come numeri reali non negativi, non come conteggi interi. Questa scelta è giustificata dal fatto che il modello è pensato per gruppi sufficientemente numerosi, per i quali le fluttuazioni sono trascurabili e l'approssimazione continua è adeguata. Inoltre, la rappresentazione continua permette di utilizzare strumenti analitici che sarebbero più complessi in un modello a conteggi interi. \par

La scelta di distinguere quattro popolazioni, invece di due o tre, è motivata dalla necessità di rappresentare separatamente:
\begin{itemize}
    \item il rischio di infortunio, che è diverso per sani, pregressi e malconci;
    \item la capacità di allenarsi, che distingue gli infortunati dagli altri;
    \item il recupero parziale, che non riporta l'atleta allo stato di sano ma lo rende nuovamente allenabile.
\end{itemize}

\subsection{Accumulo di fatica}
A queste popolazioni si aggiungono due variabili di stato che descrivono l'accumulo di fatica negli atleti:
\begin{itemize}
    \item Fatica fisica del gruppo $\bar{F}_f$: rappresenta il livello medio di affaticamento muscolare e articolare del gruppo che si allena. Aumenta con il carico di allenamento e diminuisce con il recupero.
    \item Fatica psicologica del gruppo $\bar{F}_f$: rappresenta il livello medio di affaticamento mentale del gruppo che si allena. Aumenta con il carico di allenamento e con lo stress legato alla competizione, e diminuisce con il recupero.
\end{itemize}

Le due componenti della fatica sono trattate separatamente perché evolvono con dinamiche proprie e contribuiscono in modo distinto al rischio di infortunio. Un atleta può essere fisicamente fresco ma mentalmente affaticato, o viceversa. La separazione permette di modellare scenari in cui il carico fisico e quello psicologico non sono proporzionali, e di studiare l'effetto di ciascuna componente sul rischio complessivo. Inoltre, le due componenti possono avere tassi di recupero diversi, e la loro separazione permette di catturare questa differenza senza forzare una relazione di proporzionalità tra fatica fisica e psicologica. \par

Le variabili $ \bar{F}_f(0) $ e $ \bar{F}_p(0) $ sono rappresentate come numeri reali non negativi. La loro evoluzione è descritta da equazioni che bilanciano il carico aggiunto in ogni seduta e il recupero nel tempo. In ogni passo $ t \to t+1 $, il gruppo accumula una quantità di fatica proporzionale al carico della seduta, e ne smaltisce una parte proporzionale al livello di fatica corrente.